\chapter{Deploying to the Pi}
\label{ch:deploy}

\chapterbrief
  {The machine boots into the kiosk: fullscreen, correctly rotated, GPU-accelerated, no keyboard attached.}
  {Operational discipline --- the difference between a demo and a machine.}
  {Something you can hang on a wall and walk away from.}

\section{Rotating the display}

A 7-inch panel mounted vertically needs the \emph{compositor} rotated, not the
app. Raspberry Pi OS (Bookworm and later) uses \textbf{labwc} on Wayland.

\code{\textasciitilde/.config/labwc/autostart}:

\begin{shell}
# 1024x600 panel mounted rotated -> the app sees a 600x1024 portrait screen.
wlr-randr --output HDMI-A-1 --transform 90 &
\end{shell}

Check your output name with \code{wlr-randr}; it may be \code{HDMI-A-2}.

\textbf{Touch must be rotated too}, and this is the step everyone forgets. An
unrotated touchscreen on a rotated display gives inverted axes: you press the
top-left tile and the bottom-right one activates. Maddening, and it looks like a
UI bug.

In \code{\textasciitilde/.config/labwc/rc.xml}, inside \code{<libinput>}:

\begin{xcode}
<device category="touch">
  <transformationMatrix>0 -1 1 1 0 0</transformationMatrix>
</device>
\end{xcode}

That matrix is a 90\textdegree{} rotation. For 270\textdegree{} use
\code{0 1 0 -1 0 1}. If touch ends up mirrored rather than rotated, you have the
wrong one.

\begin{concept}{Why the app does not rotate itself}
If JavaFX rendered rotated, every animation would run through a transform, the
compositor would still composite unrotated, and touch would \emph{still} be
wrong.

Rotating at the compositor means everything below it --- including input --- is
consistently transformed, and the app simply sees a portrait screen. Apply a
transformation at the lowest layer that owns all the affected data.
\end{concept}

\section{GPU acceleration: the one flag that matters}

JavaFX on a Pi will silently fall back to software rendering, and the difference
is night and day: a fade that takes 16\,ms on the GPU takes 400\,ms on the CPU,
and your kiosk feels broken.

The cause is specific. The VideoCore VI driver reports \textbf{OpenGL 2.1}, and
JavaFX's Prism pipeline checks for a higher version and refuses. The GPU is
perfectly capable; the version check is wrong.

\begin{shell}
java -Dprism.forceGPU=true \
     -Dprism.verbose=true \
     -cp "target/classes:target/dependency/*" com.midnightbrewer.ui.MainApp
\end{shell}

\code{-Dprism.verbose=true} prints the chosen pipeline at startup. You want to
see \code{es2} (hardware). If it says \code{sw}, you are on the CPU.

Also worth setting in \code{/boot/firmware/config.txt}:

\begin{shell}
dtoverlay=vc4-kms-v3d
gpu_mem=128
\end{shell}

\subsection{Other things that mattered on a Zero 2 W}

\begin{itemize}
  \item \textbf{Cache the pulse rings.} \code{ring.setCache(true);
        ring.setCacheHint(CacheHint.SPEED);} Without it, JavaFX re-rasterises the
        vector circle every frame while it scales. With it, you are scaling a
        bitmap, which is nearly free.
  \item \textbf{Hide the cursor} (\code{scene.setCursor(Cursor.NONE)}). Xwayland
        parks a pointer mid-screen; it sits on a tile and leaves it stuck in
        \code{:hover}, which reads as a pre-selected drink.
  \item \textbf{Scope hover styling} to a \code{.desktop} CSS class you enable
        only in windowed development mode. A touchscreen has no hover state to
        represent.
\end{itemize}

\section{Autostart with systemd}

\code{/etc/systemd/system/midnight-brewer.service}:

\begin{shell}
[Unit]
Description=The Midnight Brewer kiosk
After=graphical.target

[Service]
Type=simple
User=martin
Environment=WAYLAND_DISPLAY=wayland-0
Environment=XDG_RUNTIME_DIR=/run/user/1000
WorkingDirectory=/home/martin/midnight-brewer
ExecStart=/usr/bin/java -Dprism.forceGPU=true \
    -cp "target/classes:target/dependency/*" com.midnightbrewer.ui.MainApp
Restart=always
RestartSec=5

[Install]
WantedBy=graphical.target
\end{shell}

\begin{shell}
sudo systemctl daemon-reload
sudo systemctl enable --now midnight-brewer
journalctl -u midnight-brewer -f      # watch it
\end{shell}

\begin{concept}{\texttt{Restart=always} is the important line}
A kiosk that dies at three in the morning and stays dead is worse than one that
restarts in five seconds.

Combined with the broad \code{catch (RuntimeException)} from
Chapter~\ref{ch:kiosk}, you have two layers of resilience: bad taps do not kill
the worker thread, and a fatal crash does not kill the machine.

Design for the failure you cannot foresee by making recovery automatic.
\end{concept}

Remember the group membership, or Pi4J cannot open the device:

\begin{shell}
sudo usermod -aG spi,gpio martin      # log out and back in
\end{shell}

\section{The deploy loop}

Compiling on a Zero 2 W is slow enough to break your concentration. Two options.

\textbf{Rsync the source, build on the Pi} (simple, about 40 seconds):

\begin{shell}
rsync -av --exclude target/ ./ martin@raspberrypi.local:~/midnight-brewer/
ssh martin@raspberrypi.local 'cd ~/midnight-brewer && mvn -q -DskipTests compile'
\end{shell}

\textbf{Or build locally and ship the classes} (about 5 seconds). Java bytecode
is architecture-independent, so a laptop build runs unchanged on ARM:

\begin{shell}
mvn -q -DskipTests compile
rsync -av target/classes/ martin@raspberrypi.local:~/midnight-brewer/target/classes/
\end{shell}

Do the dependency copy \emph{once}:

\begin{shell}
ssh martin@raspberrypi.local \
  'cd ~/midnight-brewer && mvn dependency:copy-dependencies -DincludeScope=runtime'
\end{shell}

\begin{trap}{Stale classes}
If you change a file and the old \code{.class} survives, you will debug behaviour
that no longer exists in your source.

When something makes no sense \emph{at all} --- when the code in front of you
cannot possibly produce the output you are seeing --- run \code{mvn clean
compile} before anything else.
\end{trap}

\section{Field hygiene}

Three habits that keep a kiosk healthy.

\begin{description}[leftmargin=0pt, style=nextline]
  \item[Power-cycle the RF field per transaction.] \code{softReset()} $+$
        \code{antennaOn()} at the start of each tap puts any lingering card back
        to IDLE. Without it, a card left on the antenna after one purchase is in
        the wrong state for the next.
  \item[Always close the link.] try-with-resources on \code{Pi4jSpiLink}, and
        \code{close()} in the JavaFX \code{stop()} hook. An RF field left
        energised overnight cooks the module.
  \item[Never log key material.] It is tempting during debugging to print session
        keys. If that line survives into production, your logs are now a key
        store. Print the \emph{TI} and the \emph{counter} instead --- they tell
        you everything you need for protocol debugging and nothing an attacker
        wants.
\end{description}

\section{The finished machine}

\begin{center}
\begin{tikzpicture}[font=\footnotesize\sffamily, node distance=6mm,
   b/.style={draw=brewernavy, rounded corners=2pt, minimum width=62mm,
             minimum height=8mm, inner sep=3pt}]
  \node[b] (boot) {boot $\rightarrow$ labwc (rotated) $\rightarrow$ systemd};
  \node[b, below=of boot] (fx) {JavaFX kiosk (GPU accelerated)};
  \node[b, below=of fx, draw=brewerblue, fill=brewerblue!8] (pt) {\code{PaymentTerminal}};
  \node[b, below=of pt] (dpt) {\code{DesfirePaymentTerminal}};
  \node[b, below=of dpt] (stack) {SPI $\rightarrow$ RC522 $\rightarrow$ ISO-DEP $\rightarrow$ DESFire};
  \node[b, below=of stack, draw=brewergreen] (debit) {a real debit};
  \foreach \a/\bb in {boot/fx, fx/pt, pt/dpt, dpt/stack, stack/debit}{
    \draw[-{Stealth[length=4pt]}, brewergrey] (\a) -- (\bb);}
\end{tikzpicture}
\end{center}

\section{Commit}

\begin{shell}
git add -A
git commit -m "deploy: rotation, GPU flag, systemd unit, deploy scripts"
\end{shell}
